\documentclass{article}\usepackage[T1]{fontenc}\begin{document}
Three.flatten\_assoc identifies the two composites from strict native three-level strings to the original tagged telescope. mergeRows concatenates the rows and leaves the individual blocks intact; mapRows applies the old flattening to each row and retains the outer word. These intermediate functors need not be equal. Postcomposing each with the actual NativeBlocks.Nested.flatten gives equal functors, including on parallel arrows. Both intermediate functors preserve both span endpoints. Three.weighted\_assoc proves equality of the corresponding functors on weighted total categories for an arbitrary ordinary terminal payload. This weighting is not a comparison with the packed multiplication.

The fixed-length block proof uses common base/grade object and hom universes b,w. The global tagged and three-level results retain the earlier common native object universe u and hom universe max(u,v) for base and grades, and the original split endpoint data. Weak boundary isomorphisms are not turned into definitional equalities. The weighted law uses independent object/hom universes for the terminal payload. This is an equality of total functors: it needs no choice or normalization of a total cleavage. It does not assert an equality of independently chosen reindexing data.

Start at lines 929--953 of projects/contextual-composition/foundations/relative-construction/delivery/One\_relative\_construction\_for\_context\_content\_and\_effects.tex at 458ff36a4798e432dee5947afc402f9903a47f14. Recover two different units there: a singleton inserts a generator into the native telescope; an empty word is the unit for concatenating paths and indexes the packed free-monad unit. Read NativeBlocks.blocks, binary, Nested.word, collapse and flatten in the reviewed PR 208 tree. Read BlockCartesian.All.weighted for the total functor and WeakBoundary208Review.lean for the limit on transport to iso-comma presentations.

NativeAssoc.Path packages a fixed length with a functor into the real apex of that length, rather than identifying non-definitionally equal length sums. join\_pre and join\_last use the existing strict cuts. Removing the last primitive gives an induction principle on such diagrams, proving the two endpoint-matched products agree with the same snoc. tag\_join relates this operation to the prior NativeBlocks.binary. The universal triple of matching diagrams is tested on fixed shapes for objects and on each actual strict matching arrow for morphisms; Whole.assoc, mul\_assoc and both empty-unit laws follow. Thus associativity is not inferred from associativity of natural-number addition or from endpoint equality.

Rows.cat appends a row of variable-length blocks by reading its suffix in order. An empty suffix returns the prefix literally, a singleton suffix is pushed once, and at a longer suffix both original projections and the matching square are retained. Applying the original collapse to this concatenation gives actual multiplication of the evaluated rows, by the strict unit and associativity results. Separating the two row tags gives Rows.binary\_flatten for the full row category. Three.level then iterates exactly the zero/singleton/strict-pair presentation with those rows as letters. per\_length compares regroupings by induction on outer length; flatten\_assoc descends across tags. Functorial transport of the terminal payload gives weighted\_assoc. Both final composites preserve distinct parallel maps (mapFaithful, mergeFlattenFaithful); this does not imply fullness.

Ragged.lean uses decimal extensions and two rows with block lengths (1,0) and (1,1). The first route produces the explicit ordered four-block row; the second route has only two blocks. Their common flattened result has three primitive grades and terminal context 1234, not 1432. A genuine base arrow takes source 1 to 2 and the terminal context to 2234. The resulting native arrow is noninvertible; the weighted comparison is tested on an actual arrow between terminal finite-position payloads. A single empty row and the empty outer word also test distinct zero cases. Parallel.lean uses a free list monoid of grade endomorphisms. Reversing two composable matching squares yields an observed first component [5,1,2] rather than [1,2,5]; both final functors distinguish these composites although both endpoint arrows coincide.

The result is the fixed-span associativity law in the specified strict matching presentation, not a monad on an already constructed category of spans. Functorial action and naturality for varying span maps, identification of the weighted multiplication with PackedFree.join along the evaluation comparisons, and recovery of strict contextad algebras under the stated strict-map and normal-cleavage hypotheses remain open. No strict equality for an unrelated weak equipment presentation is claimed. Recheck the proof if the coproduct permits length-changing morphisms or the boundary model changes. Issue \#105 stays open.
\end{document}
